\section{Introduction}
\label{sec:intro}

Programmable Logic Controllers (PLCs) run the logic of safety-critical physical
infrastructure---water treatment, power distribution, manufacturing, rail and building
automation---where a control-logic defect can cause equipment damage, environmental
release, or injury. The dominant programming standard for these devices, IEC~61131-3,
and in particular its Ladder Diagram (LD) language, is therefore a natural and important
target for formal verification, and a growing body of tools now attempts it: ESBMC-PLC and
its successors, PLCverif at CERN, and model checkers such as nuXmv, among others.

Yet the field cannot easily measure its own progress. Unlike software verification for C,
where the SV-COMP benchmark suite and competition turned a collection of isolated tools into
a comparable, cumulatively improving community, PLC verification has \emph{no standard
benchmark suite}. Tool papers evaluate on private or incomparable program sets, so claims of
coverage, precision, and scalability cannot be checked against one another, and no shared
target exists against which a new technique can demonstrate improvement. The resources that
do exist fall short in one of two ways. Program corpora---such as the PLC-LD tank-control
dataset or the PPU clone-detection set---ship code but no formal properties and no expected
verdicts, so they are not verification tasks at all. And the one property-bearing suite we
are aware of covers Structured Text only, omitting graphical Ladder Diagram---the encoding
in which real programs are exchanged through PLCopen~XML and the one that dominates the
installed base. Building such a suite is not merely a matter of collecting programs: it
requires attaching, to each program, a formal property and a \emph{trustworthy} expected
verdict, and doing so in a format that any tool can run.

This paper presents that suite: \textbf{50 benchmarks over 83 program variants}, spanning
textual and graphical Ladder Diagram and Structured Text, across ten industrial domains,
each in standard PLCopen~XML or ST and paired with a formal property, a machine-checkable
expected verdict, and---for every violation---a triggering witness. The suite is packaged in
an SV-COMP--compatible layout with a schema validator and a one-command re-check harness, so
a new tool is added through an adapter rather than a reformatting of the corpus, and so the
artifact can directly seed a PLC/ICS category in SV-COMP.

The central challenge, and our main methodological contribution, is \emph{ground truth}. An
expected verdict that is silently wrong is worse than no benchmark, because it rewards
incorrect answers; and for programs of any realism the correct verdict is not obvious. One
episode from building the suite makes the point concretely. Two public logic-bomb corpora
supply pairs of a ``legitimate'' program and a ``malicious'' sibling, an apparently ready
source of safe/violation labels. The obvious safety property for these tank controllers is a
valve interlock---and asserting it would have been a mistake: reading the injected code
reveals that the malicious variants leave the valve logic \emph{untouched} and instead insert
a non-terminating loop, so the interlock holds in both and the ``obvious'' property labels
every attack \textsc{safe}. The only correct, uniform distinguishing property is
\emph{termination of the scan cycle}. The trustworthy label came from reading the mutation,
not from assuming what should be verified. We generalise this discipline into a ground-truth
methodology in which every verdict is established by one of three explicit, recorded
methods---by construction, fault injection, or audited cross-tool consensus---never by an
unexamined guess.

\medskip\noindent\textbf{Contributions.}
\begin{itemize}
  \item \textbf{A benchmark suite for IEC~61131-3 formal verification} (\S\ref{sec:suite}):
    50 benchmarks / 83 variants across textual LD, graphical LD, and ST, in ten industrial
    domains, each with a formal property, a machine-checkable expected verdict, and a witness.
    It is, to our knowledge, the first such suite to cover graphical Ladder Diagram.
  \item \textbf{A ground-truth methodology} (\S\ref{sec:groundtruth}) that triangulates every
    verdict across construction, fault injection, and audited consensus, together with the
    concrete lesson---from the non-terminating-loop episode---that trustworthy labels come
    from inspecting mutations, not from assuming the intended property.
  \item \textbf{A reusable, SV-COMP--aligned format and harness} (\S\ref{sec:principles},
    \S\ref{sec:groundtruth}): a schema-validated task layout and a one-command runner that
    re-checks every verdict, letting new tools plug in via adapters and enabling a PLC/ICS
    competition category.
  \item \textbf{Baseline results} (\S\ref{sec:eval}) for open PLC verifiers on the full
    suite, establishing reference verdicts and quantifying the input-format fragmentation
    that today prevents any single tool from running on all programs (\S\ref{sec:eval}: 43/45
    graphical variants match ground truth on ESBMC-PLC v8.4 (all 25 graphical benchmarks)).
\end{itemize}

\medskip\noindent
The complete corpus, property files, witnesses, schema, validator, and re-check harness are
archived under a citable DOI at the tool commit used for the baseline (\S\ref{sec:eval}).

\section{Related Work}
\label{sec:related}

\paragraph{Formal verification of IEC~61131-3.}
A growing set of tools verifies PLC logic. The ESBMC-PLC line~\citep{dantas_esbmc-plc_2026-3,dantas_esbmc-graphplc_2026,dantas_esbmc-plc_2026-1}
provides an SMT-based bounded-model-checking and $k$-induction frontend for Ladder Diagram and
Structured Text; PLCverif~\citep{tournier_plcverif_2022,fink_verifying_2024} at CERN targets Structured Text and drives
several back-end checkers; and general model checkers such as nuXmv~\citep{cavada_nuxmv_2014} and deductive
tools such as Why3~\citep{why3} have been applied through translation. These are the tools our
suite is built to serve. Their evaluations, however, use private or mutually incomparable
program sets with ad hoc properties, so cross-tool claims cannot be checked---precisely the gap
a shared, property-bearing suite closes. Our work is complementary rather than competing: we
provide the common target and, in \S\ref{sec:eval}, run several of these tools on it.

\paragraph{Benchmark suites and verification competitions.}
In software verification for C, SV-COMP~\citep{beyer_state_2024} and Test-Comp~\citep{testcomp} show how a
standard suite with explicit expected verdicts, machine-readable task definitions, and validated
witnesses can turn isolated tools into a comparable, cumulatively improving community. We adopt
that task model and adapt it to PLC verification. For IEC~61131-3 specifically, the closest prior
effort is the PLCOpen benchmark suite of \citet{ukegbu_benchmarks_2023}, which pairs
programs with properties but covers Structured Text only. Our suite differs in two respects: it
targets graphical Ladder Diagram---the encoding real programs are exchanged in and the one no
prior suite addresses---and it makes ground-truth provenance an explicit, recorded part of every
task (\S\ref{sec:groundtruth}) rather than an implicit assumption.

\paragraph{PLC program corpora and ICS security datasets.}
Several program collections exist but are not verification tasks. The PLC-LD tank-control
dataset~\citep{iacobelli_detection_2024} and the SWaT testbed corpus~\citep{rinieri_plc_defuser_2024} (which we ingest via the
PLC\_Defuser distribution) supply ``legitimate''/``malicious'' program pairs but no formal
properties or expected verdicts; the malicious variants realise the Ladder Logic Bomb threat
model of \citet{govil_ladder_2017}. The PPU pick-and-place clone-detection corpus~\citep{ppu}
provides graphical programs for a different task entirely (clone detection) and likewise ships no
properties. We reuse programs from all three, but our contribution is orthogonal to theirs: we
attach a formal property and a triangulated expected verdict to each, converting a program corpus
into a verification suite. This step is where the non-terminating-loop finding of
\S\ref{sec:groundtruth} arose---the ``malicious'' label alone does not tell a verifier
\emph{which} property is violated.

\paragraph{Formal semantics of IEC~61131-3.}
Our properties are stated over a per-scan execution model whose semantics rests on prior
formalisation work: the K-ST executable semantics of \citet{wang_k-st_2023}, the LD translation
rules of \citet{ebnenasir_formalizing_2023}, and the MATIEC compiler~\citep{de_sousa_matiec_2014} used by several
toolchains. The 15 feature micro-benchmarks in our suite are drawn from an independent executable
LD-semantics effort, which is why we treat their verdicts as tool-independent
(\S\ref{sec:groundtruth}). That line of work targets the \emph{meaning} of the languages; ours
targets a shared, verdict-bearing corpus for evaluating the tools that check them.

\section{Design Principles}
\label{sec:principles}

The community has no standard benchmark suite for formal verification of IEC~61131-3
Ladder Diagram programs. Existing resources fall short in one of two ways: program
corpora (e.g.\ the PLC-LD dataset, the PPU clone-detection set) ship code without formal
properties or expected verdicts, and are therefore unusable as verification tasks; and
the one property-bearing suite we are aware of (Ukegbu and Mehrpouyan) covers Structured
Text only, omitting the graphical Ladder Diagram encoding that dominates the installed
base. We therefore designed the suite from first principles, distilling six requirements
from established benchmarking practice---chiefly the SV-COMP task model of an explicit
\emph{expected verdict}, a machine-readable task definition, and a witness for every
violation---adapted to the specifics of PLC verification. Each principle is stated below
with the design decision it drove; \S\ref{sec:suite}--\ref{sec:groundtruth} report how
each is realised, and \S\ref{sec:threats} the residual risk.

\paragraph{P1. Ground truth is paramount, machine-checkable, and triangulated.}
A benchmark's entire value rests on the trustworthiness of its expected verdicts: a single
mislabelled task silently rewards a wrong answer and penalises a right one. We treat the
verdict, not the program, as the primary artifact. Every verdict is (i)~carried in a
schema-validated field rather than a comment, so it is machine-checkable and cannot drift
from the program; and (ii)~established by one of three explicit, recorded methods---by
construction, fault injection, or audited cross-tool consensus (\S\ref{sec:groundtruth})---so
that no label rests on a single unexamined judgement. This principle overrides the others:
where realism or coverage would have forced an unverifiable verdict, we excluded the task.

\paragraph{P2. A verification task is a program \emph{and} a formal property.}
What distinguishes a verification suite from a program corpus is that each task states,
formally, \emph{what is to be proved}. Every benchmark therefore pairs its program with one
or more properties drawn from a small, closed vocabulary (invariant, mutual-exclusion,
reachability, termination), expressed over program variables under the per-scan execution
model. Fixing a closed vocabulary keeps tasks comparable across tools and prevents the
property language from silently becoming a second, unspecified dimension of difficulty.

\paragraph{P3. Diversity across language, domain, and difficulty.}
A suite that is uniform measures one narrow capability. We require spread along three axes.
\emph{Language:} all three practically-used IEC~61131-3 forms---textual and graphical Ladder
Diagram and Structured Text---with deliberate emphasis on graphical LD, the encoding real
programs are exchanged in and the one no prior suite targets. \emph{Domain:} ten industrial
sectors, so that results are not an artefact of a single process type. \emph{Difficulty:} a
gradient from single-construct micro-benchmarks to multi-PLC controllers, so the suite can
separate tools rather than saturate at one level.

\paragraph{P4. Discrimination over triviality.}
A suite on which every tool scores perfectly is uninformative. We therefore balance the
expected verdicts (roughly half \textsc{safe}, half \textsc{violation}), so that a strategy
of always answering ``safe'' fails half the tasks, and we retain a hard tier whose members
current tools are not guaranteed to solve. Balanced verdicts also force tools to demonstrate
both proof and refutation, and expose the \emph{unknown} outcomes that a safe-only suite
would hide.

\paragraph{P5. Realistic provenance, with the authored share disclosed.}
Realism and control are in tension: real programs are representative but come without
properties or clean fault labels, while authored programs are controllable but risk tuning
the suite to the authors' own tool. We resolve this by preferring third-party and prior
programs wherever they exist, authoring only to fill genuine gaps (domains and constructs with
no public program), and disclosing the exact split (\S\ref{sec:suite}) so a reader can weigh
the two. Authored programs follow standard control idioms---interlocks, guarded actuation---rather
than any tool-specific pattern.

\paragraph{P6. A reusable, tool-agnostic format and one-command re-check.}
A suite becomes a standard only if others can run it. We adopt a stable, documented,
schema-validated layout aligned with the SV-COMP task model (task descriptor, property file,
per-violation witness), so a new tool is added through an adapter rather than a reformatting
of the corpus, and so the suite can seed an SV-COMP PLC/ICS category directly. Reproducibility
is built in: a single harness re-runs every task and checks the result against its recorded
verdict, and where a tool cannot ingest a program we record that explicitly rather than
dropping the task, keeping coverage honest.

\section{The Benchmark Suite}
\label{sec:suite}

The suite comprises \textbf{50 benchmarks} over \textbf{83 program variants},
each an IEC~61131-3 program in standard PLCopen~XML (Ladder Diagram) or Structured
Text, paired with a machine-checkable property file and an expected verdict. A
\emph{benchmark} is a verification task; a \emph{variant} is one program file under
that task (30~benchmarks bundle a safe program with one or more faulted counterparts,
\S\ref{sec:groundtruth}). Every benchmark validates against the suite's JSON-Schema
and is packaged in an SV-COMP--compatible layout (task descriptor, property file,
witness fields), so the artifact can seed a future SV-COMP PLC/ICS category without
rework.

\subsection{Composition}
Table~\ref{tab:composition} gives the breakdown by IEC~61131-3 language and by
industrial domain. The language split (15~textual~LD, 25~graphical~LD, 10~ST)
deliberately favours \emph{graphical} LD: it is the encoding in which real,
third-party programs are actually exchanged, and no prior formal-verification suite
targets it (the closest, the PLCOpen suite of Ukegbu and Mehrpouyan, covers 40~ST
programs only). The suite spans \textbf{ten industrial domains}; water treatment is
the largest (\S\ref{sec:threats}) because the two richest public ICS corpora---SWaT
and the PLC-LD tank-control dataset---are both water-domain.

\begin{table}[t]
  \centering\small
  \caption{Suite composition by language and by domain (50 benchmarks).}
  \label{tab:composition}
  \begin{tabular}{@{}llr@{\hspace{2em}}lr@{}}
    \toprule
    \multicolumn{2}{@{}l}{\textbf{Language}} & \textbf{\#} &
    \textbf{Domain} & \textbf{\#}\\
    \midrule
    Ladder Diagram & (textual)   & 15 & water treatment    & 16\\
    Ladder Diagram & (graphical) & 25 & manufacturing      & 6\\
    Structured Text&             & 10 & motor control      & 5\\
    \cmidrule(r){1-3}
    \multicolumn{2}{@{}l}{\textbf{Total}} & \textbf{50} & packaging & 5\\
                   &             &    & traffic            & 4\\
    \multicolumn{2}{@{}l}{\textbf{Difficulty}} &  & building automation & 3\\
    easy           &             & 26 & elevator           & 3\\
    medium         &             & 18 & HVAC               & 3\\
    hard           &             &  6 & power substation   & 3\\
                   &             &    & chemical batch     & 2\\
    \bottomrule
  \end{tabular}
\end{table}

\subsection{Provenance}
To balance realism against controlled coverage, the suite draws from four external
or prior sources and supplements them with authored programs
(Table~\ref{tab:provenance}). Nineteen benchmarks come from third-party or prior
corpora: the PLC-LD tank-control dataset (Iacobelli et al.), the SWaT testbed via the
PLC\_Defuser dataset, five curated ESBMC-PLC security scenarios, and one program from
the PPU pick-and-place clone-detection corpus. The remaining 31 are authored: 15 are
feature micro-benchmarks that systematically exercise individual LD constructs
(contacts, coils, seal-in, TON/TOF/TP timers, CTU/CTD counters, edge detectors,
latches), and 16 are domain control programs written to give each of the ten domains a
genuine, non-analogical anchor---particularly the domains for which no public program
was available (chemical batch, power substation).

\begin{table}[t]
  \centering\small
  \caption{Provenance of the 50 benchmarks.}
  \label{tab:provenance}
  \begin{tabular}{@{}lr l@{}}
    \toprule
    \textbf{Source} & \textbf{\#} & \textbf{Nature}\\
    \midrule
    PLC-LD dataset (Iacobelli et al.) & 8  & graphical LD, real third-party\\
    SWaT / PLC\_Defuser               & 6  & multi-PLC ST, real testbed\\
    ESBMC-PLC security scenarios      & 4  & curated ST\\
    PPU clone-detection corpus        & 1  & graphical LD\\
    Authored: feature micro-benchmarks& 15 & one LD construct each\\
    Authored: domain control programs & 16 & interlocks + 5 syntax twins\\
    \midrule
    \textbf{Total}                    & \textbf{50} & \\
    \bottomrule
  \end{tabular}
\end{table}

\subsection{Properties and coverage}
Each benchmark carries one or more formal properties in a YAML schema with a closed
vocabulary of kinds: \emph{invariant} (29 properties), \emph{termination}
(14), \emph{mutual\_exclusion} (6), and \emph{reachability} (3, used for
counterexample/trigger synthesis). Properties are expressed over program variables
under the per-scan execution model. Beyond safety invariants and interlocks, the suite
covers a property class absent from existing PLC benchmarks: \emph{termination} of the
scan cycle, which is the safety-relevant manifestation of the non-terminating-loop
logic bombs described in \S\ref{sec:groundtruth}.

\subsection{Syntax-coverage pairs}
Five benchmarks are graphical-LD renderings of a textual-LD or ST original with
identical logic (recorded via a \texttt{syntax\_pair\_of} field). These are not
padding: they let the suite act as a \emph{differential} test of whether a tool that
accepts one PLCopen encoding produces the same verdict on the other, a discrepancy we
consider a defect. They are counted once toward each language slice and flagged
explicitly, so a reader can include or exclude them.

\section{Ground-Truth Methodology}
\label{sec:groundtruth}

A benchmark suite is only as useful as its expected verdicts: a single mislabelled
task silently rewards wrong answers. We therefore establish every verdict by one of
three methods and record the method, per variant, in the benchmark descriptor.
Across the 83 variants the expected verdicts are 49~\textsc{safe} and
34~\textsc{violation}; the three methods account for 36, 29, and 18 variants
respectively (Table~\ref{tab:groundtruth}).

\begin{table}[t]
  \centering\small
  \caption{Ground-truth method by variant (83 variants).}
  \label{tab:groundtruth}
  \begin{tabular}{@{}lr l@{}}
    \toprule
    \textbf{Method} & \textbf{\#} & \textbf{Basis of the verdict}\\
    \midrule
    Expert / by-construction     & 36 & property holds (or fails) by the program's design\\
    Fault injection              & 29 & mutant of a safe base with a documented defect\\
    Cross-tool consensus         & 18 & $\geq$2 tools agree, plus manual audit\\
    \bottomrule
  \end{tabular}
\end{table}

\paragraph{Expert / by-construction.}
The feature micro-benchmarks and the authored domain interlocks are constructed so
that their property is decided by inspection: a series interlock makes two outputs
mutually exclusive; a fill-valve rung is guarded by a not-full contact. For these we
record the verdict as \textsc{expert} and give the one-line justification in the
property file.

\paragraph{Fault injection.}
For 29 benchmarks the \textsc{violation} variant is a \emph{mutant} of a known-safe
base, with the seeded defect and a triggering witness documented in the descriptor.
This is the most trustworthy source of \textsc{violation} labels because the defect is
known by construction. Crucially, both public corpora we ingest already supply such
pairs: every ``malicious'' program in the PLC-LD dataset and in SWaT is a
fault-injected sibling of a ``legitimate'' one.

\paragraph{The non-terminating-loop logic bomb.}
Inspecting these mutants surfaced a uniform, and initially non-obvious, fact: the
injected defect in \emph{every} malicious variant of both corpora is a
\emph{non-terminating loop}---a \texttt{WHILE} whose induction variable never reaches
its bound (e.g.\ \texttt{i := i*i} from zero, or a counter that is never advanced),
gated on a specific sensor value. The controllers' functional logic (valve
interlocks, pump hysteresis) is left untouched. A na\"ive labelling that asserted a
valve-safety property would therefore have marked these mutants \textsc{safe}, because
that property still holds---the label would have been wrong. The correct, uniform
distinguishing property is \emph{termination of the scan cycle}: the safe program
completes every scan, the mutant hangs once the trigger input occurs. We encode this
as the \texttt{termination} property kind, which a bounded model checker discharges via
a scan-cycle watchdog. This episode is the strongest argument for the methodology:
the trustworthy label came only from reading the injected code, not from assuming what
the ``obvious'' safety property should be.

\paragraph{Cross-tool consensus.}
For the \textsc{safe} variants of the fault-injection pairs, where no by-construction
argument fixes the verdict, we require agreement between independent tools followed by
manual audit of any disagreement, and record the confirming tools in the descriptor.
Disagreements are not discarded: they are adjudicated by hand and reported, since a
program on which mature tools disagree is among the most informative in the suite.

\paragraph{Reproducibility.}
Verdicts are re-checkable end to end: \texttt{run\_all.sh} runs every variant through
the verifier---\texttt{-{}-k-induction} for expected-\textsc{safe} tasks,
\texttt{-{}-incremental-bmc} for expected-\textsc{violation} tasks, with the scan
watchdog enabled---and compares the result to the recorded verdict, emitting a
per-variant pass/fail table. The full suite, property files, witnesses, and this
harness are archived with a citable DOI.

\section{Evaluation}
\label{sec:eval}

We establish reference verdicts on the suite with ESBMC-PLC v8.4, built from source with the
Ladder Diagram front-end enabled (\texttt{-DENABLE\_LD\_FRONTEND=On}); expected-\textsc{safe}
tasks are run under $k$-induction and expected-\textsc{violation} tasks under incremental BMC,
both with the scan-cycle watchdog enabled. The harness attempts all 50 benchmarks and records
the outcome of every variant against its recorded verdict.

\paragraph{Results.}
Of the 83 variants, the 45 that the front-end accepts---\textbf{all 25 graphical PLCopen~XML
benchmarks}---run, and \textbf{43 match the recorded verdict with zero disagreements}; the two
non-matches are \textsc{unknown} outcomes ($k$-induction proof-strength results on safe
edge-detector programs, reported as \textsc{unknown} rather than folded into \textsc{safe}, per
\S\ref{sec:threats}). This confirms the recorded ground truth across every authored domain
interlock (safe variant proved, faulted variant refuted with a counterexample) and every
non-termination tank controller (the injected logic bomb caught by the scan watchdog). It is
direct, tool-checked evidence that the property files, the seeded defects, and the witnesses are
consistent. (Reaching this required aligning the property syntax to v8.4's C-style boolean
operators and running an adapter that maps the suite schema to the tool's property format---the
adapter is released with the suite.)

\paragraph{A format-and-semantics finding.}
The remaining 38 variants---the textual-Ladder and Structured-Text benchmarks---did not run, and
the reason is more than skin-deep. First, format: v8.4's front-end accepts \emph{only}
PLCopen~XML, so a program in the textual LD form or a stand-alone ST module is rejected outright.
Second, and more interesting, \emph{semantics}: when we transcribed the textual timer and counter
benchmarks into PLCopen~XML, their recorded verdicts---established under an independent executable
LD semantics---did \emph{not} transfer, because ESBMC-PLC's scan-cycle models of \texttt{TON}/
\texttt{CTU} and of set/reset latches differ from that reference. Porting these benchmarks
therefore requires re-establishing ground truth under each tool's timing model, not a mechanical
re-serialisation. Both are concrete instances of the fragmentation this suite is designed to
expose---one tool's front-end admits a single serialisation, and even the meaning of a standard
timer is not agreed---and both motivate the suite carrying every program in vendor-neutral
PLCopen~XML \emph{with per-tool verdicts where the semantics diverge}, which is the principal item
of remaining work.

\paragraph{Scope of these numbers.}
These verdicts establish an ESBMC-PLC reference column, not yet a cross-tool comparison.
Running PLCverif and nuXmv on the same programs requires their respective input formats (SCL
and SMV), reinforcing the same fragmentation point; that multi-tool baseline, and the
PLCopen~XML normalisation that enables it, are the immediate next steps.

\section{Threats to Validity}
\label{sec:threats}

We discuss threats along the standard construct, internal, external, and
reproducibility axes, and state the mitigations we adopted and the residual risk.

\paragraph{Construct validity: do the tasks measure verification capability?}
The properties are written by us, so the suite risks measuring ``can the tool prove
\emph{the properties we chose}'' rather than program correctness. We mitigate this by
deriving each property from the program's own control logic---interlocks, guarded
actuation, scan termination---rather than inventing abstract assertions, and by drawing
from a closed, documented vocabulary (\S\ref{sec:suite}). Most benchmarks carry a single
property; the suite is not intended to certify full functional correctness, only to
provide decidable, well-motivated safety obligations. A second construct threat is
\emph{fault-model coverage}. The 34 \textsc{violation} variants comprise 14 scan
non-terminations, 14 safety-invariant violations, and 6 interlock (mutual-exclusion)
violations. Importantly, the two public logic-bomb corpora we ingest (PLC-LD and SWaT)
inject \emph{only} non-terminating loops; every other defect class in the suite
(interlock bypass, actuator manipulation, overflow-guard evasion) is represented through
\emph{authored or curated} mutants. Real-world violation examples in the suite are thus
narrower in defect type than the authored ones, a limitation inherited from the available
public corpora.

\paragraph{Internal validity: are the expected verdicts correct?}
A mislabelled task is the most damaging defect a suite can have. Our three-method
triangulation (\S\ref{sec:groundtruth}) targets this directly: fault-injection labels are
known by construction, expert labels are decidable by inspection, and consensus labels
require independent agreement plus manual audit. The non-terminating-loop episode
(\S\ref{sec:groundtruth}) is evidence the audit process catches mislabelling that a
plausible-but-wrong property would have introduced. Two residual risks remain. First,
the expert and fault-injection labels rely on our reading of the programs; we mitigate by
publishing every justification and witness so they are independently checkable, and by the
end-to-end re-check below. Second, a \emph{consensus} \textsc{safe} label and the
tool under evaluation are not fully independent: ESBMC-PLC+ participates both in
establishing consensus and in the baseline. We limit this by requiring at least one
tool-independent basis (by-construction or a second, architecturally different checker)
for every \textsc{safe} label, and by noting that only 18 of 83 variants rest on
consensus at all. We further distinguish \emph{proof strength} from \emph{truth}: a
$k$-induction run may return \emph{unknown} on a genuinely safe program; such outcomes are
reported as \textsc{unknown}, not silently folded into \textsc{safe}.

\paragraph{External validity: do results generalise?}
Two facts bound generalisation. First, 31 of 50 benchmarks are authored, which risks a
suite tuned to constructs our own frontend handles well; we counter this by including 19
third-party/prior programs, by sourcing the authored feature benchmarks from an
independent executable-semantics effort, and by writing the authored domain programs to
standard interlock patterns rather than to our tool. Second, the domain distribution is
skewed: water treatment accounts for 16 of 50 benchmarks, because the two richest public
ICS corpora are both water-domain. We report all results per domain so a reader can
reweight, and we ensured every one of the ten domains has at least two genuine anchors so
none is represented only by analogy. Program scale is a further external threat: many
benchmarks are small, and heavy scalability stress rests on the multi-PLC SWaT programs
and the counter-fuse scalability family alone. Finally, five benchmarks are
syntax-coverage twins; they are flagged and counted once per language slice so their shared
logic cannot inflate coverage claims.

\paragraph{Reproducibility.}
All verdicts are re-checkable with a single command (\texttt{run\_all.sh},
\S\ref{sec:groundtruth}), which selects the proof mode per expected verdict and compares
the tool's result to the recorded label. Two practical caveats attach to cross-tool
baselining: the baseline checkers consume different input formats, so no single tool runs
on all 50 programs, and we report per-tool coverage explicitly rather than dropping
unsupported rows---the format fragmentation is itself a finding that motivates a common
suite. The complete corpus, property files, witnesses, schema, and harness are archived
under a citable DOI at the exact tool commit used for the baseline.

\section{Conclusion}
\label{sec:conclusion}

Formal verification of IEC~61131-3 programs has active tools but no common yardstick: without
a shared, property-bearing benchmark suite, progress cannot be measured and techniques cannot be
compared. We have presented such a suite---50 benchmarks over 83 program variants, spanning
textual and graphical Ladder Diagram and Structured Text across ten industrial domains, each
carrying a formal property, a machine-checkable expected verdict, and, for every violation, a
triggering witness. To our knowledge it is the first formal-verification benchmark suite to cover
graphical Ladder Diagram, the encoding in which real programs are exchanged. The corpus is
packaged in an SV-COMP--compatible layout with a schema validator and a one-command re-check
harness, so a new tool is added through an adapter and the artifact can directly seed a PLC/ICS
category in SV-COMP.

Our central lesson concerns ground truth. For programs of any realism the correct verdict is not
self-evident, and an authoritative-looking label can be quietly wrong: the ``malicious'' variants
of two public corpora would have been mislabelled \textsc{safe} by the obvious valve-interlock
property, because the injected defect is a non-terminating loop that leaves the valve logic
intact. The trustworthy label came only from reading the mutation. We therefore make ground-truth
provenance a first-class, recorded property of every task, established by construction, fault
injection, or audited cross-tool consensus---a discipline we believe any verification benchmark,
in any domain, should adopt.

The suite is designed to grow. Natural next steps are broadening the fault model beyond the
non-termination bombs that dominate the available public corpora, adding programs exported from
proprietary vendor toolchains to widen manufacturer coverage, and enlarging the corpus toward the
scale needed for a standing competition category. By releasing the benchmarks, property files,
witnesses, schema, and harness under a citable DOI, we intend the suite to serve as a shared,
extensible foundation on which the PLC formal-verification community can measure and compare its
progress.

